% Solutions of 28-sequential-circuits.
\section*{Chapter \ref{ch:b1:sequential-circuits} --- Sequential Circuits and Finite-State Machines}

\begin{solution}{exo:b1:sequential-circuits:1}
$Q = 1$ (set), $1$ (hold), $0$ (reset), $1$ (set), $1$ (hold).
\end{solution}

\begin{solution}{exo:b1:sequential-circuits:2}
6, 7, 0, 1, 2. From 5 ($101$) bits 0 and 1 flip; from 6 bit 0; from 7 all three; from 0 bit 0; from 1 bits 0 and 1. Bit
$i$ flips when all the bits below it are 1.
\end{solution}

\begin{solution}{exo:b1:sequential-circuits:3}
The \omterm{def:b1:sequential-circuits:fsm}{states} are $\varepsilon/0$, then $1/0$, $1/0$, $10/0$, $1/1$, $10/0$, $1/1$, $1/0$, $10/0$, and the output word
$00010100$: the detector fires at the fourth and the sixth bit, the second $101$ overlapping the first.
\end{solution}

\begin{solution}{exo:b1:sequential-circuits:4}
$1101$, $1110$, $0111$, $0011$: each bit enters on the left and the rightmost one leaves.
\end{solution}

\begin{solution}{exo:b1:sequential-circuits:5}
Twice. Once rightly, while the count is 2; once wrongly, between 3 and 4, where the \omterm{def:b1:sequential-circuits:register}{counter} shows 2 for $d$ time units
(\cref{out:b1:sequential-circuits:counter}). Two repairs: use the synchronous \omterm{def:b1:sequential-circuits:register}{counter}, whose bits change together; or feed
the comparator's output into a \omterm{def:b1:sequential-circuits:flip-flop}{flip-flop} on the same clock, so that the bell sees it only at the edges, after the value
has settled.
\end{solution}

\begin{solution}{exo:b1:sequential-circuits:6}
The pair \code{("red", 1)} has no \omterm{def:b1:sequential-circuits:fsm}{transition}: the \omterm{def:b1:sequential-circuits:moore-mealy}{state diagram} leaves out what happens when someone presses the button
while people cross, and the first word never does that. A \omterm{def:b1:sequential-circuits:fsm}{finite-state machine} must have a \omterm{def:b1:sequential-circuits:fsm}{transition} for every state and
every input symbol. Add \code{("red", 1): "green"} (the press is ignored, the light goes on), or keep the light red one
more period; either is a decision, and the test should check that the table has $3 \times 2$ entries.
\end{solution}

\begin{solution}{exo:b1:sequential-circuits:7}
\omcode{code/bachelor-1/28-sequential-circuits/exercises/detect_0110.py}{12}{19}{The detector, checked against the
definition}
The detector of \cref{lst:b1:sequential-circuits:lock}, built for $0110$, has the \omterm{def:b1:sequential-circuits:fsm}{states} $e$, $0$, $01$, $011$; after
$0110$ it goes to $0$, since the last 0 begins a new pattern. It agrees with the definition on all 8191 words.
\end{solution}

\begin{solution}{exo:b1:sequential-circuits:8}
\omcode{code/bachelor-1/28-sequential-circuits/exercises/mod6.py}{11}{14}{The next-state equations}
The program finds $A^+ = BC + AC'$, $B^+ = BC' + A'B'C$, $C^+ = C'$: 9 gates. From 5 ($101$): $A^+ = 0 + 0 = 0$, $B^+ = 0 + 0 =
0$, $C^+ = 0$, back to 0; without the don't-cares the equations are longer. The run prints $0, 1, \dots, 5, 0, \dots$.
\end{solution}

\begin{solution}{exo:b1:sequential-circuits:9}
$\lceil \log_2 k \rceil$ with a binary encoding, $k$ with one-hot. The \omterm{def:b1:sequential-circuits:datapath}{control unit} has 23 \omterm{def:b1:sequential-circuits:fsm}{states}: 5 \omterm{def:b1:sequential-circuits:flip-flop}{flip-flops} against 23.
One-hot buys simple logic: the output signals are ORs of the \omterm{def:b1:sequential-circuits:fsm}{states} that order them, with no decoding of the state, and
the next-state input of each \omterm{def:b1:sequential-circuits:flip-flop}{flip-flop} is an OR of the \omterm{def:b1:sequential-circuits:fsm}{transitions} entering its state. Hardware designers often pay the
\omterm{def:b1:sequential-circuits:flip-flop}{flip-flops} for that speed.
\end{solution}

\begin{solution}{exo:b1:sequential-circuits:10}
Mealy $M$ to Moore $M'$: let $q_i$ be the state of $M$ after $a_1 \dots a_i$ and $o_i = \lambda(q_{i-1}, a_i)$ its $i$-th
output. Claim: $M'$ is in $(q_i, o_i)$ after $i$ symbols ($i \geq 1$). For $i = 1$, $\delta'((q_0, 0), a_1) = (\delta(q_0,
a_1), \lambda(q_0, a_1)) = (q_1, o_1)$. If it holds for $i$, then $\delta'((q_i, o_i), a_{i+1}) = (q_{i+1}, o_{i+1})$. So the
$i$-th output of $M'$, the output of the state entered, is $o_i$. Moore $M$ to Mealy $M'$: both machines go through the
same \omterm{def:b1:sequential-circuits:fsm}{states}, by induction since they have the same $\delta$; the $i$-th output of $M'$ is $\lambda(\delta(q_{i-1}, a_i)) =
\lambda(q_i)$, the $i$-th output of $M$. In both directions the output is not delayed: the Mealy output on the \omterm{def:b1:sequential-circuits:fsm}{transition}
into $q_i$ and the Moore output of $q_i$ are read after the same symbol.
\end{solution}

\begin{solution}{exo:b1:sequential-circuits:11}
$T \geq 50 + 12 \times 25 + 30 = 380$ ps, so at most $1 / 380\,\text{ps} \approx 2.63$ GHz. The earliest change at a
\omterm{def:b1:sequential-circuits:flip-flop}{flip-flop} input comes $50 + 2 \times 25 = 100$ ps after the edge, well after the 10 ps \omterm{def:b1:sequential-circuits:setup}{hold time}. With no gate it comes
after 50 ps, still enough: the clock-to-output delay itself protects the \omterm{def:b1:sequential-circuits:setup}{hold time}, unless the clock reaches the two
\omterm{def:b1:sequential-circuits:flip-flop}{flip-flops} at different times.
\end{solution}

\begin{solution}{exo:b1:sequential-circuits:12}
\omcode{code/bachelor-1/28-sequential-circuits/exercises/vending.py}{13}{16}{The \omterm{def:b1:sequential-circuits:moore-mealy}{state diagram}, as a function}
The credit is 0, 1 or 2 ($AB = 00, 01, 10$); with coins $c_1 c_0 = 00, 01, 10$, the code 11 and the credit 3 are
don't-cares. The program finds $A^+ = Bc_0 + A'B'c_1 + Ac_1'c_0'$, $B^+ = Bc_1'c_0' + A'B'c_0$, $v = Bc_1 + Ac_0 + Ac_1$, $r =
Ac_1$: 16 gates, and the circuit agrees with the machine on all 9841 sequences of up to eight coins.
\end{solution}

\begin{solution}{pb:b1:sequential-circuits:1}
\textbf{1.} \omterm{def:b1:sequential-circuits:fsm}{States} $e, 1, 10, 101$. From $e$: $0 \to e$, $1 \to 1$. From $1$: $0 \to 10$, $1 \to 1$. From $10$: $0 \to e$,
$1 \to 101$. From $101$: $0 \to 10$, $1/\mathbf{1} \to 1$. Every other output is 0.

\textbf{2.} After $1011$, the longest end of the input that begins the pattern is the last $1$: a new attempt has begun.
$e$ would forget it, and miss the lock opening on $1011011$.

\textbf{3.} Line 22: it tries the ends of $w$ from the longest to the empty one, and keeps the first that is a prefix.

\textbf{4.} The outputs are $0001001$: the lock opens after the fourth and the seventh keys, the second $1011$ sharing a 1
with the first.

\textbf{5.} The double loop over prefixes and keys makes one entry per pair, and \code{next} always finds an end, the
empty one at worst.

\textbf{6.} As the fourth key is pressed: the output is 1 during the clock cycle of that key, before the next edge.

\textbf{7.} $e/0, 1/0, 10/0, 101/0$ and the new $1/1$, the state $1$ entered by the opening \omterm{def:b1:sequential-circuits:fsm}{transition}.

\textbf{8.} Only reachable pairs are kept: only $1$ is ever entered with an output of 1.

\textbf{9.} Both print $0001001$, as \cref{prop:b1:sequential-circuits:conversion} says; the \omterm{def:b1:sequential-circuits:moore-mealy}{Moore machine} shows each 1 as
the output of the state $1/1$.

\textbf{10.} $1/1$ is the only state with output 1, so its own output already differs from every other's.

\textbf{11.} If two \omterm{def:b1:sequential-circuits:fsm}{states} differ on a word $w$, two inputs leading to them (here the prefixes that name them) followed by
$w$ give different outputs; a machine in the same state after both would give the same.

\textbf{12.} The printed words separate all $\binom{4}{2} = 6$ Mealy pairs and all $\binom{5}{2} = 10$ Moore pairs, so no
two \omterm{def:b1:sequential-circuits:fsm}{states} can be merged: in a \omterm{def:b1:sequential-circuits:moore-mealy}{Moore machine}, after the words $1$ and $1011$ the futures are equal but the current
outputs differ.

\textbf{13.} Mealy: $e = 00$, $1 = 01$, $10 = 10$, $101 = 11$. Moore: $e/0 = 000$, $1/0 = 001$, $10/0 = 010$, $101/0 = 011$,
$1/1 = 100$.

\textbf{14.} From $10$ ($AB = 10$) on 1 to $101$ ($11$): $B^+ = 1$. From $101$ on 0 to $10$: $B^+ = 0$. The low bit
records whether the last key was 1.

\textbf{15.} No state has them, so their rows never occur: don't-cares.

\textbf{16.} Only $1/1$ has $A = 1$, and the codes 5 to 7 are don't-cares.

\textbf{17.} Mealy: NOT $x$, NOT $B$, ANDs $Bx'$, $AB'x$, $ABx$, one OR: 6. Moore: NOTs of $A$, $B$, $C$, $x$, five ANDs
and two ORs: 11.

\textbf{18.} The Mealy output is one AND gate away from the key, within the same cycle; the Moore output comes from a
\omterm{def:b1:sequential-circuits:flip-flop}{flip-flop}, one clock edge later but free of the input's glitches. Timing depends on the clock as much as on gates.

\textbf{19.} $2^{13} - 1 = 8191$, the empty word included.

\textbf{20.} By induction on the length: they start in the same state, and if they are in the same state after $i$ keys,
equal next-state functions put them in the same state after $i + 1$, with the same output.

\textbf{21.} Every state followed by every key appears within the first five keys of some word, so the check covers every
reachable (state, key) pair, and \textbf{20} extends it to every length.

\textbf{22.} Two \omterm{def:b1:sequential-circuits:fsm}{states} with one code would merge them, and some word in the check, such as the distinguishing word
following the two prefixes, would show a difference.

\textbf{23.} Keys 2 to 9 send every state back to $e$ (they begin no attempt); the \omterm{def:b1:sequential-circuits:fsm}{states} stay the same. There are
$(10^{13} - 1) / 9 \approx 1.1 \times 10^{12}$ words of up to 12 keys; \textbf{20} and \textbf{21} make a check of all \omterm{def:b1:sequential-circuits:fsm}{states}
and keys enough.

\textbf{24.} \emph{Mealy saves a state}: the lock needs 4 Mealy \omterm{def:b1:sequential-circuits:fsm}{states} and 5 Moore \omterm{def:b1:sequential-circuits:fsm}{states}, each number minimal by the
distinguishing words; the circuits need 2 \omterm{def:b1:sequential-circuits:flip-flop}{flip-flops} and 6 gates against 3 \omterm{def:b1:sequential-circuits:flip-flop}{flip-flops} and 11 gates.
\end{solution}
